% chapters/00_page3_assignment.tex
% Trang 3 - Nhiệm vụ luận văn thạc sĩ
% Theo mẫu biểu mẫu 11
% Author: Nguyễn Tấn Khôi
% Date: 2025-12-14

%% ============================================================================
%% TRANG 3: NHIỆM VỤ LUẬN VĂN THẠC SĨ
%% ============================================================================

\thispagestyle{empty}

% Header với 2 cột
\noindent
\begin{tabular}{@{}p{0.48\textwidth}p{0.48\textwidth}@{}}
\fontsize{12}{14}\selectfont
ĐẠI HỌC QUỐC GIA TP.HCM\newline
\textbf{TRƯỜNG ĐẠI HỌC BÁCH KHOA}
&
\fontsize{12}{14}\selectfont
\mbox{\textbf{CỘNG HÒA XÃ HỘI CHỦ NGHĨA VIỆT NAM}}\newline
\textit{Độc lập - Tự do - Hạnh phúc}
\end{tabular}

\vspace{1cm}

\begin{center}
{\fontsize{16}{19}\selectfont\bfseries
NHIỆM VỤ LUẬN VĂN THẠC SĨ\\
}
\end{center}

\vspace{0.5cm}

{\fontsize{12}{14}\selectfont

\noindent
\begin{tabular}{@{}p{9.5cm}p{3.5cm}@{}}
Họ tên học viên: NGUYỄN TẤN KHÔI & MSHV: \studentid\\
Ngày, tháng, năm sinh: 08/02/1999 & Nơi sinh: TP.HCM\\
Chuyên ngành: \major & Mã số: \majorcode\\
\end{tabular}

\vspace{0.8cm}

\textbf{I. TÊN ĐỀ TÀI:} \mbox{\thesistitleVN} - \thesistitleEN\\


\textbf{II. NHIỆM VỤ VÀ NỘI DUNG:}
\vspace{0pt}

\begin{itemize}[leftmargin=0.5cm, itemsep=0pt, topsep=0pt]
\item Nghiên cứu tổng quan về các phương pháp tránh va chạm trong hệ đa robot.
\item Tìm hiểu thuật toán Proximal Policy Optimization (PPO) và các kỹ thuật huấn luyện multi-agent.
\item Đề xuất cải tiến về learning rate scheduling, value clipping và chiến lược huấn luyện hai giai đoạn.
\item Thiết kế và chế tạo robot di động sử dụng cảm biến LiDAR và Jetson Nano.
\item Triển khai thuật toán lên robot thực tế và đánh giá.
\end{itemize}

\textbf{III. NGÀY GIAO NHIỆM VỤ:} (Ghi theo trong QĐ giao đề tài)\\

\textbf{IV. NGÀY HOÀN THÀNH NHIỆM VỤ:} (Ghi theo trong QĐ giao đề tài)\\

\textbf{V. CÁN BỘ HƯỚNG DẪN} (Ghi rõ học hàm, học vị, họ, tên): \advisorname\\

}


\hfill \textit{Tp. HCM, ngày . . . . tháng . . . . năm 2025}


\noindent
\begin{tabular}{@{}p{7.5cm}p{7.5cm}@{}}
\centering\textbf{CÁN BỘ HƯỚNG DẪN} & \centering\textbf{CHỦ NHIỆM BỘ MÔN ĐÀO TẠO}\tabularnewline
\centering\textit{(Họ tên và chữ ký)} & \centering\textit{(Họ tên và chữ ký)}\tabularnewline
\end{tabular}

\vspace{2.0cm}

\begin{center}
\textbf{TRƯỞNG KHOA ĐIỆN -- ĐIỆN TỬ}\\
\textit{(Họ tên và chữ ký)}
\end{center}

\clearpage
